\scanpage{68}
% archive foliation: 67
\deleted{Un épinglage de $M$ s'identifie donc à un}

On peut considérer, indépendamment de la question d'épingler $\underline J$, \deleted{celle de trouver} \inserted{l'épinglage} ``relatif'' \deleted{une section de} consistant en \deleted{une} la donnée d'une section du torseur $P_M$ sous $\underline{\operatorname{End}}(M_0)$, ce qui équivaut \inserted{donc} aussi à la donnée d'un morphisme
\[
\tag{47}\underline J\longrightarrow\deleted{\unsure{V_{\underline J}(}}M\deleted{\unsure{)}}
\]
qui relève le morphisme d'inclusion
\[
\underline J\hookrightarrow M_0
\]
et qui soit une section de $V_{\underline J}(M)$ i.e. qui soit à trace nulle (rel. à $S$).

On a en fait automatiquement un morphisme
\[
\underline J\longrightarrow M\setminus M_0
\]
donc en composant avec \deleted{\unsure{l'inclusion}}
\[
\diagram{\begin{tikzcd}[ampersand replacement=\&]
(\underline J\arrow[r,dotted]\& )M\arrow[r,hook]\&E\arrow[r]\&\Sigma_{\underline J}\simeq E/(1,\sigma_E)
\end{tikzcd}}
\]
on trouve un morphisme
\[
\tag{48}\underline J\longrightarrow\Delta_t
\]
et la condition de relèvement pour (47) signifie aussi que le composé de (48)
